\chapter{Hadamard matrices from finite fields}\label{ch:paley}


A row of four signs can be orthogonal to another row of four signs, even though none of its entries is zero. For example, $(1,1,1,1)$ and $(1,-1,1,-1)$ have dot product $0$. Can we arrange $n$ rows of $n$ signs so that every pair is orthogonal?

Such an arrangement is called a Hadamard matrix. Its defining equation, $HH^{\mathsf T}=nI$, says that every row has squared length $n$ and distinct rows meet at right angles. The first construction takes one small matrix and doubles it. The second builds signs from the squares and non-squares of a finite field.

The finite-field construction brings three earlier ideas together: cyclic multiplicative groups, quadratic characters, and the sum of a character over a field. The formulas look technical until we notice that orthogonality is just a count of matching and mismatching signs.

\section{Sylvester's doubling}

The \emph{Kronecker product} of an $m\times m$ matrix $A$ and an $n\times n$ matrix $B$ is the $mn\times mn$ block matrix $A\otimes B=(a_{ij}B)$. It
satisfies $(A\otimes B)(C\otimes D)=(AC)\otimes(BD)$ and $(A\otimes B)^{\mathsf T}=A^{\mathsf T}\otimes B^{\mathsf T}$, by multiplying blocks.

\begin{proposition}[Sylvester, 1867]\label{conjectures:prop:kron}
If $H$ and $K$ are Hadamard matrices of orders $m$ and $n$, then $H\otimes K$ is a Hadamard matrix of order $mn$. In particular the matrices
$S_1=(1)$, $S_{2^{k+1}}=\begin{pmatrix}S_{2^k}&S_{2^k}\\S_{2^k}&-S_{2^k}\end{pmatrix}=\begin{pmatrix}1&1\\1&-1\end{pmatrix}\otimes S_{2^k}$ are
Hadamard matrices of order $2^k$.
\end{proposition}
\begin{proof}
$(H\otimes K)(H\otimes K)^{\mathsf T}=(HH^{\mathsf T})\otimes(KK^{\mathsf T})=(mI_m)\otimes(nI_n)=mnI_{mn}$, and the entries are products of signs.
\end{proof}

The rows of $S_{2^k}$ are the \emph{Walsh functions}: the entry in row $i$ and column $j$ is $(-1)^{i\cdot j}$, where $i\cdot j$ is the dot product of
the binary expansions. They are the characters of the group $(\mathbb Z/2)^k$, and multiplying by $S_{2^k}$ is the
Walsh--Hadamard transform of signal processing.

\section{Paley's constructions}

Let $q$ be an odd prime power and $\mathbb F_q$ the field with $q$ elements (Theorem~\ref{thm:fq}). Its multiplicative group is cyclic (Theorem~\ref{thm:cyclicmult}); let $g$ be a generator. The \emph{quadratic character} $\chi:\mathbb F_q\to\{0,\pm1\}$ is $\chi(0)=0$, $\chi(x)=1$ if $x$ is a nonzero
square and $\chi(x)=-1$ otherwise.

\begin{lemma}\label{conjectures:lem:chi}
\begin{enumerate}
\item $\chi(g^j)=(-1)^j$, so $\chi$ is multiplicative and there are $\frac{q-1}2$ nonzero squares.
\item $\chi(-1)=(-1)^{(q-1)/2}$: $-1$ is a square if and only if $q\equiv1\pmod4$.
\item $\sum_{x\in\mathbb F_q}\chi(x)=0$.
\item For $c\neq0$, $\sum_{x\in\mathbb F_q}\chi(x)\chi(x+c)=-1$.
\end{enumerate}
\end{lemma}
\begin{proof}
(a) The squares of $g^i$ are the $g^{2i}$, so the nonzero squares are the even powers of $g$, which are $\frac{q-1}2$ elements since $g$ has even order
$q-1$. Hence $\chi(g^j)=(-1)^j$, and $\chi(g^ig^j)=(-1)^{i+j}=\chi(g^i)\chi(g^j)$; with $0$ the multiplicativity is clear. (b) $-1$ is the element of
order $2$, that is $g^{(q-1)/2}$. (c) There are as many $+1$ as $-1$. (d) The term $x=0$ is $0$. For $x\neq0$,
$\chi(x)\chi(x+c)=\chi(x)^2\chi(1+cx^{-1})=\chi(1+cx^{-1})$. As $x$ runs over $\mathbb F_q^*$, $1+cx^{-1}$ runs over $\mathbb F_q\setminus\{1\}$, so the sum is
$\sum_y\chi(y)-\chi(1)=0-1$.
\end{proof}

Index rows and columns by the elements of $\mathbb F_q$, in any order, and let $Q$ be the \emph{Jacobsthal matrix} $Q_{xy}=\chi(x-y)$. Write $J$ for an
all-ones matrix and $\mathbf j$ for an all-ones column, with sizes given by the context.

\begin{lemma}\label{conjectures:lem:jacobsthal}
$Q\mathbf j=0$, $\mathbf j^{\mathsf T}Q=0$, $QQ^{\mathsf T}=qI-J$, and $Q^{\mathsf T}=\chi(-1)Q$: the matrix $Q$ is symmetric if $q\equiv1\pmod4$ and
skew-symmetric if $q\equiv3\pmod4$.
\end{lemma}
\begin{proof}
The row and column sums are $\sum_y\chi(x-y)=0$ and $\sum_x\chi(x-y)=0$, by Lemma~\ref{conjectures:lem:chi}(c). The entry $(x,z)$ of $QQ^{\mathsf T}$ is
$\sum_y\chi(x-y)\chi(z-y)$; put $u=x-y$: it is $\sum_u\chi(u)\chi(u+z-x)$, which is $q-1$ if $z=x$ and $-1$ otherwise, by (a) and (d). Finally
$Q_{yx}=\chi(y-x)=\chi(-1)\chi(x-y)$.
\end{proof}

\begin{theorem}[Paley I, 1933]\label{conjectures:thm:paley1}
Let $q\equiv3\pmod4$ be a prime power and
\[
S=\begin{pmatrix}0&\mathbf j^{\mathsf T}\\-\mathbf j&Q\end{pmatrix}.
\]
Then $H=I+S$ is a Hadamard matrix of order $q+1$.
\end{theorem}
\begin{proof}
The entries of $S$ are $0$ on the diagonal and $\pm1$ elsewhere, because $\chi(x-y)=\pm1$ for $x\neq y$; so $H$ is a $\pm1$ matrix. Since $Q$ is
skew-symmetric, so is $S$. By Lemma~\ref{conjectures:lem:jacobsthal},
\[
SS^{\mathsf T}=\begin{pmatrix}0&\mathbf j^{\mathsf T}\\-\mathbf j&Q\end{pmatrix}\begin{pmatrix}0&-\mathbf j^{\mathsf T}\\\mathbf j&Q^{\mathsf T}\end{pmatrix}
=\begin{pmatrix}\mathbf j^{\mathsf T}\mathbf j&\mathbf j^{\mathsf T}Q^{\mathsf T}\\Q\mathbf j&\mathbf j\mathbf j^{\mathsf T}+QQ^{\mathsf T}\end{pmatrix}
=\begin{pmatrix}q&0\\0&J+qI-J\end{pmatrix}=qI.
\]
So $HH^{\mathsf T}=(I+S)(I-S)=I-S^2=I+SS^{\mathsf T}=(q+1)I$.
\end{proof}

\begin{theorem}[Paley II, 1933]\label{conjectures:thm:paley2}
Let $q\equiv1\pmod4$ be a prime power and
\[
S=\begin{pmatrix}0&\mathbf j^{\mathsf T}\\\mathbf j&Q\end{pmatrix},\qquad H=\begin{pmatrix}S+I&S-I\\S-I&-S-I\end{pmatrix}.
\]
Then $H$ is a Hadamard matrix of order $2(q+1)$.
\end{theorem}
\begin{proof}
Now $Q$ and $S$ are symmetric, and the computation of the previous proof, with $+\mathbf j$ in place of $-\mathbf j$, gives $S^2=SS^{\mathsf T}=qI$. The
entries of $S\pm I$ are $\pm1$, and $H$ is symmetric, so
\[
\begin{aligned}HH^{\mathsf T}=H^2&=\begin{pmatrix}(S+I)^2+(S-I)^2&(S+I)(S-I)-(S-I)(S+I)\\(S-I)(S+I)-(S+I)(S-I)&(S-I)^2+(S+I)^2\end{pmatrix}\\
&=\begin{pmatrix}2(q+1)I&0\\0&2(q+1)I\end{pmatrix},\end{aligned}
\]
since $S$ commutes with $I$ and $(S+I)^2+(S-I)^2=2S^2+2I$.
\end{proof}
